\section{Reverse probes and fixed reservoirs}

\label{c:probes}

Fix a set of observed cards and leave $h\ge n/8$ labels hidden,
starting from any deterministic placement. Run two walkers independently
conditional on the observed paths, from any two available positions.
After five sweeps their collision probability, averaged over the observed
paths, is at least $(1-n^{-1/100})/h$. Thus the mean collision matrix
contains almost the full uniform matrix as an entrywise lower bound.
The proof follows possible partner positions backward to the initial
hidden set, using reverse paths close to independent uniform probes.
The resulting matrix bound implies a squared-norm contraction; iteration
then produces simultaneous endpoint caps for eight fixed reservoirs.

\begin{theorem}[Reverse-probe route]\label{c:probes:main}
For the physical Thorp shuffle on $n=2^d$ positions,
\[
 \|q_d^{*(100000d+1)}-U_{S_n}\|_{\mathrm{TV}}\longrightarrow0.
\]
The pairwise collision estimate used in its proof is uniform over the
initial positions of every fixed observed set.
\end{theorem}

We first prove the collision estimate, then construct one retained
probability law with eight simultaneous reservoir caps, and finally
apply the eight-block representation estimate. The additional physical
shuffle at the end of the schedule supplies exact parity cancellation.

\label{c:probes:p001}

Write \(n=N=2^d\), and encode the positions (numbered from \(0\) at the top) by words of length
\(d\) in binary, starting with the most significant bit. A step taking time \(k-1\) to time
\(k\) sends the two positions \(0x,1x\) to \(x0,x1\), where \(x\) has length \(d-1\), by a fair
random bijection. Call this a switch, indexed by the time and \(x\); the switches use
independent coins. In particular the shuffle acts by permutations of the positions, independent
of the card labels. We use this description throughout.

\label{c:probes:p003}
\label{c:probes:p004}

\subsection{A conditional walk estimate}

\label{c:probes:p005}

Track the full paths of specified cards during a specified interval,
starting from known positions. The positions occupied by the remaining
cards are available; let their number be \(h\). By the common flow of
Lemma~\ref{lem:marginal-averaging}, in the physical coordinates used
here, the conditional transition matrix \(K\) has rows indexed by
initial available positions and columns by terminal available positions.
A switch with one available input has its available output prescribed;
with two available inputs, either output bit is chosen with probability
\(1/2\). The matrix describes a start at any available position without
requiring the occupant's identity. It is doubly stochastic by the step
rule, or by averaging bijections of the available sets. The new task is
to compare two walks that are independent conditional on the observed
paths; the following construction keeps that conditional independence
explicit.

\label{c:probes:p006}

\subsubsection{Estimating availability by reverse paths}

When two conditional walkers occupy the same position, they separate
with probability $1/2$ if the other input of their switch is available,
and otherwise stay together. We first estimate these successive
availability tests. The estimate conditions on one or two prescribed
paths of an ordinary independent-switch process, called the base
process. It does not condition on the observed-card paths used to
define $K$.

\begin{lemma}[Reverse probes under prescribed base paths]
\label{lem:c-reverse-probes}
Set $B=5d$, $L=\lfloor d/2\rfloor$, and $E_d=100d^3/n$.
Fix integers $0\le a\le u$, $u-a\ge d$, $0\le v\le L$, and
$u+v\le B$. Fix the entire base history through time $a$, including
an available set $A_a$ of size $h$, and let the later available sets be
its images under the base process. Put $p=h/n$.

Condition further on a feasible specified base path $W$ from $a$
through $u+v$. Index its bits by their times of appending, writing
$\alpha_k$ for the bit appended at time $k$, with the same indexing
for the starting word. One may also specify a second base path from
$a$ through $u$, provided it agrees with $\alpha$ on bit indices
$u-d+1,\ldots,u-d+v$. No other future information is conditioned on.
For $u\le s<u+v$, let $Z_s$ be obtained from $W_s$ by
complementing its first bit. Then, under this conditioning,
\begin{equation}
 \left|\mathbb E\!\left[2^{-\#\{u\le s<u+v:Z_s\in A_s\}}\right]
             -(1-p/2)^v\right|\le E_d.
 \label{c:probes:eq:3}
\end{equation}
\end{lemma}

\begin{proof}
\label{c:probes:p010}
By the path-cylinder argument in Lemma~\ref{lem:marginal-averaging},
the prescribed paths fix their visited switches and leave all other
base coins after $a$ independent and fair.
From each $Z_s$, trace backward through the base process to time $a$;
membership in $A_s$ is equivalent to the trace's endpoint lying in
$A_a$. Compare these traces with independent ideal probes, starting
at the same words and prepending independent fair bits while dropping
the last bit at each reverse step.

\label{c:probes:p011}
A reverse step from $k$ to $k-1$ uses the switch indexed by the
length-$(d-1)$ prefix at time $k$. Expose reverse steps in decreasing
time order. Until two ideal probes request the same switch, or a probe
requests a prescribed-path switch, assign the requested base coin its
ideal value. These are distinct unconstrained switches, so the
assignments have the required independent fair law. Complete every
other unconstrained base coin ordinarily. Thus the true and ideal
traces agree unless such a conflict occurs.

\label{c:probes:p012}
Put $\ell_s=s-d+1$. The probe starting at $Z_s$ differs from
$\alpha$ exactly at index $\ell_s$ in its starting word. At a reverse
step with $k\ge\ell_s+1$, this bit remains in its prefix and prevents
a conflict with $W$. It also prevents a conflict with the optional
second path in that path's time range, by the agreement assumption.
For two active probes, use the larger starting time $s$. The earlier
probe has the unchanged bit $\alpha_{\ell_s}$: their starting times
differ by less than $L<d$, so that index belongs to both starting
words and is not the bit complemented by the earlier probe. This again
prevents a conflict while $k\ge\ell_s+1$.

When $k\le\ell_s$, the entire prefix of the indicated ideal probe
consists of $d-1$ fresh reverse bits. Its probability of matching a
fixed prescribed prefix, or the prefix of an independent ideal probe,
is $2^{-(d-1)}$. These are bounds for the unrestricted ideal probes;
we do not condition on previous avoidance of conflicts. A union bound
over at most $5d$ times and $2L+\binom L2$ comparisons per time gives
conflict probability at most $E_d$.

Every probe has run at least $d$ steps by time $a$. Its endpoint is
therefore uniform on $V$, independently of the other probes, and has
probability $p$ of belonging to the fixed set $A_a$. The ideal
expectation in \eqref{c:probes:eq:3} is consequently $(1-p/2)^v$.
The tested variable lies in $[0,1]$, so the coupling error is at most
$E_d$. The case $v=0$ is immediate.
\end{proof}

\subsubsection{From survival to a collision lower bound}

\begin{proposition}[A fixed-reservoir collision bound]
\label{c:probes:kernel-contraction}
Let $h\ge n/8$ and observe the paths of a fixed set of $n-h$ cards
through $B=5d$ physical shuffles, starting from arbitrary deterministic
positions. Let $K$ be the resulting conditional transition matrix on
the available positions, and put $Q=\mathbb E[KK^{\mathsf T}]$.
For all sufficiently large $d$, every pair $i,j$ of initial available
positions satisfies
\begin{equation}
 Q(i,j)=\mathbb E\sum_x K(i,x)K(j,x)
       \ge\frac{1-n^{-1/100}}h,
 \label{c:probes:collision-minorization}
\end{equation}
where the sum is over the terminal available positions and the
expectation is over the observed paths. Consequently, for every real
row vector $z$ on the initial available positions with sum zero,
\begin{equation}
\mathbb E\|zK\|_2^2\ \le\ n^{-1/100}\|z\|_2^2
\label{c:probes:eq:1}
\end{equation}
for all sufficiently large \(d\), uniformly over the observed set and its initial placement.
The vector norms are Euclidean.
\end{proposition}

Lemma~\ref{noise:halves} already gives the squared-norm contraction,
with the smaller factor $(31/32)^{3d}+32e^{-n/256}$ when
$\alpha=1/8$ and $t=5d$. The additional conclusion here is
\eqref{c:probes:collision-minorization}: it gives a lower bound of order
$1/h$ for each pair of starting positions, rather than only a bound on
the quadratic form on sum-zero vectors.

\begin{proof}
Put $p=h/n$ and call the initial time zero.

\label{c:probes:p007}

The matrix $Q$ is symmetric and stochastic; its entry \(Q(i,j)\) is the probability that two walkers
starting from available positions \(i,j\), run independently given the tracked paths, are at the same
position at time \(B\). The diagonal entries are at least \(1/h\), since the sum of squares of
a probability vector of length \(h\) is at least \(1/h\).

We will prove $Q(i,j)\ge(1-n^{-1/100})/h$ for every pair $i,j$.
Reverse probes estimate how long the two conditional walkers remain
together. We then force their last meeting and sum over its possible
times. The entrywise lower bound finally gives contraction on the
sum-zero vectors.

\label{c:probes:p008}

We use the following joint experiment for the two walkers, named \(W,W'\). Run the base
switch process generating, in particular, the tracked paths and
moving the available set. Also sample an extra independent fair bit at each time. Let \(W\)
always follow the base process from its position. At switches different from the switch of
\(W\), let \(W'\) also follow the base process from its position. When both use the same switch
and that switch has two available inputs, use that time's extra bit as the appended bit for \(W'\). If
they use the same switch and it has only one available input (so they are at the same position), \(W'\)
also follows the base process, since the available output is forced. This keeps both walkers in
available positions and gives the required law. Indeed, conditional on the tracked paths, the joint
transition at each step given the two walkers' histories is the product of the prescribed
conditional transitions: the base coins on switches with two available inputs are independent conditional
on the tracked paths, as explained above (and the walkers' histories so far only use coins at
earlier times); the extra coin also gives independence when both walkers have two choices at the same
switch.

\label{c:probes:p009}

In this experiment, without conditioning on future tracked paths, whenever the walkers are
distinct, their two appended bits are independent and fair given the past. This follows from
independence of the step's switches when using different switches (even when there is only one
available input at such a switch), and from the extra independent bit when using the same switch from
distinct positions. Also, if they are at the same position and we consider the chance of
staying together over an interval of steps, the relevant probabilities given the full base
process are 1 or \(1/2\) at each such step. Specifically, at a time \(s\) (for the step to
\(s+1\)), denote by \(Z_s\) the position obtained from the position of \(W\) by complementing
its first bit. It is the other input of the switch. If \(Z_s\) is available, an independent extra
coin must give agreement in order to stay together; otherwise they will agree automatically.
Thus, starting together, the probability given the base process of staying together, using
fresh extra coins over steps from times \(s=u,\dots,u+v-1\), is
\begin{equation}
2^{-\#\{s=u,\ldots,u+v-1: Z_s\text{ is available at time }s\}}.                 \label{c:probes:eq:2}
\end{equation}

Use the constants from Lemma~\ref{lem:c-reverse-probes}:
\[
L=\lfloor d/2\rfloor,\qquad E_d=100 d^3/n.
\]
The survival probability is therefore approximately $(1-p/2)^v$
whenever the specified-path hypotheses hold. Summing the possible last
meeting times will produce the factor
$\sum_{r\ge0}(1-p/2)^r=2/p$; paired with a meeting cost $1/(2n)$,
this gives the desired collision scale $1/h$.

\label{c:probes:p013}

For the pair experiment, a newly meeting pair at some time \(m'\ge d+1\) (distinct at \(m'-1\),
equal at \(m'\)) must have appended equal bits for the last \(d\) steps. They were distinct
immediately before each of those steps, since any separation after being together would put
unequal bits in that string of steps. Consequently the probability of newly meeting at that
time is at most \(2^{-d}\), by the conditional fair bit property for a distinct pair.

\label{c:probes:p014}

We now turn the probe estimate into a collision lower bound by forcing a last meeting near the
terminal time. Take \(r=0,\ldots,L-1\), and put
\[
m=B-r,\qquad b=m-d-1.
\]
We will use that
\begin{equation}
\mathbb P(W_b=W'_b)\ \le\ (1-p/2)^L+E_d+ L/n.                     \label{c:probes:eq:4}
\end{equation}
Here \(W_s,W'_s\) denote their positions at time \(s\). If together at \(b\), they have either
newly met in the preceding \(L\) steps (each such meeting at a time at least \(d+1\)), or were
together continuously from \(b-L\). For the latter event, put
$u=b-L$. Given the full base process and the walker experiment through
$u$, its probability is the indicator of $W_u=W'_u$ times the
survival probability in \eqref{c:probes:eq:2}, with $v=L$.
Drop that indicator before averaging: retaining it could bias the
availability tests. Lemma~\ref{lem:c-reverse-probes} now applies with
$a=u-d\ge0$ and only the specified base path of $W$. Its upper bound
$(1-p/2)^L+E_d$, together with the $L/n$ bound for new meetings, proves
\eqref{c:probes:eq:4}.

\label{c:probes:p015}

Given the past at $b$ with the walkers distinct, put $c=m-d$ and
require their appended bits to differ at $c$ and agree at
$c+1,\ldots,m$. The differing bit keeps the positions distinct through
time $m-1$, so this requirement has probability exactly
$2^{-(d+1)}=1/(2n)$ and forces a new meeting at $m$. We claim that,
given the paths through this meeting, their probability of staying
together for the remaining $r$ steps is at least
\begin{equation}
(1-p/2)^r-E_d.                                                     \label{c:probes:eq:5}
\end{equation}
To justify the conditioning, let $\widetilde W$ be the ordinary base
path starting from $W'_b$. Through time $m-1$ the walker $W'$ follows
$\widetilde W$. At the first step, if it shares $W$'s switch, the
required opposite output is its own base output. At subsequent steps
before $m$ the bit at index $c$ prevents a shared switch. At $m$ the
two base paths use the same switch, but $W'$ takes $W$'s output by its
extra coin. Thus their appended bits have the following form:
\[
\begin{array}{c|ccc}
\text{bit index}&W&W'&\widetilde W\\ \hline
c&\alpha_c&1-\alpha_c&1-\alpha_c\\
k=c+1,\ldots,m-1&\alpha_k&\alpha_k&\alpha_k\\
m&\alpha_m&\alpha_m&1-\alpha_m
\end{array}
\]
Now condition on the particular paths in this table and on $W$ through
time $B$. This specifies exactly the base coins along $W$ and
$\widetilde W$ through their indicated times, together with the extra
coins used at a shared switch. These assignments suffice to force the
meeting event: both base paths start at available positions, remain
available, and have two available inputs whenever they share a switch.
There is no additional restriction on future availability. All other
base coins and future extra coins remain independent and fair.

The optional second path in Lemma~\ref{lem:c-reverse-probes} is
$\widetilde W$, not the walker after meeting. With $a=b$, $u=m$,
and $v=r$, its required agreement interval is $c+1,\ldots,c+r$.
It agrees there by the table, since $r<L<d$; its opposite final bit
at $m$ is outside this interval. Averaging the survival probability
\eqref{c:probes:eq:2} under precisely this conditioning therefore gives
\eqref{c:probes:eq:5}.

\label{c:probes:p016}

The events requiring the indicated meeting and then staying together are disjoint for different
\(m\) (each requires being distinct just before that meeting). Using \eqref{c:probes:eq:4} and
\eqref{c:probes:eq:5}, and writing \(q=1-p/2\) and \(f=q^L+E_d+L/n\), for each \(r\) the event
contributes at least \(\big((1-f)q^r-E_d\big)/(2n)\) to the equality probability: the chance of
being distinct at \(b\) is at least \(1-f\), the bit requirement given the past and
distinctness there has probability \(1/(2n)\), and the conditional chance of staying then is at
least \(q^r-E_d\). We obtain for any starts \(i,j\)
\[
\begin{aligned}
Q(i,j) &\ge \frac{1}{2n}\left((1-f)\sum_{r=0}^{L-1}q^r - L E_d\right)\\
 &=\frac{1}{h}\left((1-f)(1-q^L)-\frac{p}{2} L E_d\right)
 \ \ge\ \frac{1}{h}(1-n^{-1/100}),
\end{aligned}
\]
where the last bound holds for all sufficiently large \(d\), since \(q\le 15/16\), so \(q^L\le
(15/16)^{d/2-1}=o(n^{-1/100})\), and the other error terms are bounded by polynomials in \(d\)
times \(1/n\). With \(\delta=n^{-1/100}\), \(Q\) can therefore be written as \(1-\delta\) times
the uniform stochastic matrix plus \(\delta\) times another symmetric stochastic matrix. The
latter has Euclidean operator norm at most 1 (by Jensen's inequality and the column sums). Thus
\(z Q z^{\mathsf T}\le\delta\|z\|_2^2\) for zero-sum \(z\), giving \eqref{c:probes:eq:1}.
\end{proof}

\label{c:probes:p017}

\subsection{Obtaining a measure with several spread-out projections}

\label{c:probes:p018}

Now suppose that fixed cards are tracked from deterministic initial positions over \(t_0=2000
B\) steps, still with \(h\ge n/8\). For a particular untracked card, let \(w\) be its
conditional distribution on the terminal available positions. By iterating \eqref{c:probes:eq:1},
\begin{equation}
\mathbb E\|w-\mathbf u_h\|_2^2\le n^{-20},                              \label{c:probes:eq:6}
\end{equation}
where \(\mathbf u_h\) is the vector assigning mass \(1/h\) to each available position there. In detail, when
we observe the tracked paths block by block, the distribution of this card given the
observations is successively multiplied by the corresponding \(K\). Indeed the new tracked
paths give no further information on the occupant identities at available positions at the start of that
block, since those paths use new switches and do not depend on these identities. Conditional on
observations so far the next block is as in \eqref{c:probes:eq:1}. The uniform vector is
carried to the uniform vector by each matrix; \eqref{c:probes:eq:1} now applies to the
difference at each block, starting with squared norm at most 1. In particular the probability
that \(\|w-\mathbf u_h\|_2>n^{-3}\) is at most \(n^{-14}\).

\label{c:probes:p019}

Identify the cards initially by their positions and write \(G=\mathfrak S_n\) for the group
acting on positions (so \(g h\) applies \(h\) first). Partition the positions into eight
disjoint sets of size \(n/8\), for \(d\ge3\), and let \(H_j\) (\(1\le j\le8\)) be the subgroup
permuting the \(j\)-th set only. Denote by \(\mu\) the distribution on \(G\) of the permutation
from \(t_0\) successive steps, taking initial positions to terminal positions. We will have a
probability measure \(\nu\) on \(G\) with \(\|\mu-\nu\|_{\mathrm{TV}}=o(1)\) such that for each
\(j\),
\begin{equation}
\nu(gH_j)\le 2\,\frac{|H_j|}{|G|}\qquad(g\in G).                   \label{c:probes:eq:7}
\end{equation}

\label{c:probes:p020}

To construct it, for each set use a fixed ordered list of the complementary cards. Consider
each card on that list given the full paths of the preceding cards on the list. Its conditional
distribution at the end satisfies \eqref{c:probes:eq:6}, with \(h\) the number of available positions when
only those preceding cards are tracked. Require \(\|w-\mathbf u_h\|_2\le n^{-3}\) for all these
considerations on the eight lists. By a union bound, the resulting good event on the base
process with deterministic initial cards has probability tending to 1. Use the law of the
permutation conditioned on this event as \(\nu\); its total variation distance from \(\mu\) is
at most the probability of the complement event, by writing \(\mu\) as the mixture. For a
particular list to have specified distinct endpoints, and all requirements for that list to
hold, the probability is at most
\[
\prod_{i=1}^{n-n/8}\frac{1+n^{-2}}{n-i+1}
\]
by successive conditioning on paths of preceding cards: the requirement at index \(i\) is
determined by these paths, and when it holds the probability of the specified endpoint for the
card at index \(i\) is at most \(1/(n-i+1)+n^{-3}\) (or zero if it is not available). Formally
this product bound follows by taking conditional expectations starting with the last index.
Specifying the endpoints of that entire list gives the coset \(gH_j\). The good event implies
all requirements for the list; dividing the probability bound by the probability of the good
event gives \eqref{c:probes:eq:7} for large \(d\), since \(\prod_{i=1}^{n-n/8}
(n-i+1)^{-1}=|H_j|/|G|\) and \((1+n^{-2})^{n-n/8}\to1\).

\label{c:probes:p021}

\subsection{The eight-coset application}
\label{c:probes:eight-coset-application}
The representation estimate needed here has the following exact form.
Suppose $H_1,\ldots,H_8$ permute disjoint equal blocks of an $n$-point set,
and a probability $\nu$ on $S_n$ satisfies
$\nu(gH_j)\le2|H_j|/|S_n|$ for every $g,j$. For every irreducible
$\rho_\lambda$ of dimension $D_\lambda$ other than the trivial and sign
representations, its restriction to the product of the $H_j$ has
$T_\lambda$ distinct tensor types, and the eight-coset transfer gives
\begin{equation}\label{c:probes:eq:8}
 \|\widehat\nu(\lambda)\|_{\mathrm{op}}
 \le\sqrt{2T_\lambda}\,D_\lambda^{-3/8}.
\end{equation}
The accompanying type count gives $2T_\lambda\le D_\lambda^{1/4}$
uniformly for sufficiently large $n$, and hence
$\|\widehat\nu(\lambda)\|_{\mathrm{op}}\le D_\lambda^{-1/4}$.
These operator estimates are Theorem~\ref*{l-l:joint-type-operator}
and its eight-block specialization,
Corollary~\ref*{l-l:eight-block-probe} of~\cite{partialFourier}.
Separately, the paragraph following that corollary supplies the direct
summation argument
\begin{equation}\label{c:probes:eq:9}
 \sum_{\lambda\ne(n),(1^n)}D_\lambda^{-3}=o(1),
\end{equation}
using the inverse-third-power majorant. The orbit-span estimate
includes all multiplicity spaces; its hypothesis is a cap on the
right-multiplication cosets $gH_j$.
Thus \eqref{c:probes:eq:7} has exactly the required orientation and
normalization.

\label{c:probes:p029}

\label{c:probes:bound}
Compose ten independent permutations with law \(\nu\), and then apply one actual step of the
shuffle, independently; call the distribution \(\eta\). The sign has expectation zero due to
the last step (one fair switch suffices to balance it), and the last step always has Fourier
operator norm at most 1. Plancherel now gives the following bound for squared \(L^2\) deviation
of the density from uniform:
\[
\begin{aligned}
\mathbb E_{g\sim \text{unif}(G)} (|G|\eta(g)-1)^2
 &= \sum_{\pi\ne \text{trivial}} D_\pi \|\widehat\eta(\pi)\|_{\rm HS}^2
 \ \le\ \sum_{\pi\notin\{\text{trivial, sign}\}} D_\pi^2\,(D_\pi^{-1/4})^{20}=o(1).
\end{aligned}
\]
Here we used the product rule by independence and bounded the Hilbert-Schmidt norm by the
square root of the dimension times the operator norm. Total variation deviation tends to zero
by Cauchy-Schwarz. Replacing \(\nu\) by \(\mu\) in these compositions changes the law by at
most \(10\|\mu-\nu\|_{\rm TV}=o(1)\) in total variation, by coupling or contraction under
independent composition. With \(\mu\) these are exactly independent shuffle increments of
\(t_0\) steps each on the positions, followed by the extra shuffle step. Since the starting
deck can be identified with any labeling of the initial positions, this is a worst-case bound.
It proves Theorem~\ref{c:probes:main}, and gives \(t_{\rm mix}(d)\le100000d+1\) for all
sufficiently large \(d\).

\label{c:probes:p030}

Together with the support obstruction \eqref{mart:lower-exact}, this
route gives $2d-O(1)\le t_{\mathrm{mix}}(d)\le100000d+1$ for all
sufficiently large $d$. The reverse-probe argument also retains the
entrywise collision bound \eqref{c:probes:collision-minorization} for
every fixed reservoir and every pair of available starting positions.
